%
\documentclass[12pt]{article}

\usepackage[usenames]{color}
\usepackage{verbatim}
\usepackage{fancyhdr}
\usepackage{graphicx}
\usepackage{amssymb,amsthm}
\usepackage{bm}
%\usepackage[fleqn]{amsmath}
\usepackage{amsmath}
% \usepackage[final,notref,notcite,color]{showkeys}
\usepackage[left=1in,top=1in,right=1in,bottom=1in,letterpaper]{geometry}
\renewcommand{\baselinestretch}{1.2}

\pagestyle{fancy}
\lhead{}
\chead{}
\rhead{}
\lfoot{}
\cfoot{\thepage}
\rfoot{}



%% editing
\usepackage[shortlabels]{enumitem}
\usepackage[normalem]{ulem}

%% macros for commenting
\usepackage{mathtools}
\usepackage[normalem]{ulem} % to use \sout

%% math packages
% \usepackage{amsmath} % not needed for Springer
\usepackage{amssymb}
\usepackage{mathrsfs}

%\usepackage{subfigure}
%\usepackage{url}

%% a light-weight algorithm environment
\newtheorem{algo}{Algorithm}

%% highlighting and commenting
\newcommand{\outline}[1]{{\color{brown}#1}}
\newcommand{\rev}[1]{{\color{blue}#1}}
\newcommand{\remove}[1]{{\sout{#1}}}
\newcommand{\cut}[1]{{}}

%% macros for letters

\newcommand{\va}{{\mathbf{a}}}
\newcommand{\vb}{{\mathbf{b}}}
\newcommand{\vc}{{\mathbf{c}}}
\newcommand{\vd}{{\mathbf{d}}}
\newcommand{\ve}{{\mathbf{e}}}
\newcommand{\vf}{{\mathbf{f}}}
\newcommand{\vg}{{\mathbf{g}}}
\newcommand{\vh}{{\mathbf{h}}}
\newcommand{\vi}{{\mathbf{i}}}
\newcommand{\vj}{{\mathbf{j}}}
\newcommand{\vk}{{\mathbf{k}}}
\newcommand{\vl}{{\mathbf{l}}}
\newcommand{\vm}{{\mathbf{m}}}
\newcommand{\vn}{{\mathbf{n}}}
\newcommand{\vo}{{\mathbf{o}}}
\newcommand{\vp}{{\mathbf{p}}}
\newcommand{\vq}{{\mathbf{q}}}
\newcommand{\vr}{{\mathbf{r}}}
\newcommand{\vs}{{\mathbf{s}}}
\newcommand{\vt}{{\mathbf{t}}}
\newcommand{\vu}{{\mathbf{u}}}
\newcommand{\vv}{{\mathbf{v}}}
\newcommand{\vw}{{\mathbf{w}}}
\newcommand{\vx}{{\mathbf{x}}}
\newcommand{\vy}{{\mathbf{y}}}
\newcommand{\vz}{{\mathbf{z}}}

\newcommand{\ta}{{\tilde{a}}}
\newcommand{\tb}{{\tilde{b}}}
\newcommand{\tc}{{\tilde{c}}}
\newcommand{\td}{{\tilde{d}}}
\newcommand{\te}{{\tilde{e}}}
\newcommand{\tf}{{\tilde{f}}}
\newcommand{\tg}{{\tilde{g}}}

\newcommand{\ti}{{\tilde{i}}}
\newcommand{\tj}{{\tilde{j}}}
\newcommand{\tk}{{\tilde{k}}}
\newcommand{\tl}{{\tilde{l}}}
\newcommand{\tm}{{\tilde{m}}}
\newcommand{\tn}{{\tilde{n}}}

\newcommand{\tp}{{\tilde{p}}}
\newcommand{\tq}{{\tilde{q}}}

\newcommand{\ts}{{\tilde{s}}}

\newcommand{\tu}{{\tilde{u}}}
\newcommand{\tv}{{\tilde{v}}}
\newcommand{\tw}{{\tilde{w}}}
\newcommand{\tx}{{\tilde{x}}}
\newcommand{\ty}{{\tilde{y}}}
\newcommand{\tz}{{\tilde{z}}}

\newcommand{\vA}{{\mathbf{A}}}
\newcommand{\vB}{{\mathbf{B}}}
\newcommand{\vC}{{\mathbf{C}}}
\newcommand{\vD}{{\mathbf{D}}}
\newcommand{\vE}{{\mathbf{E}}}
\newcommand{\vF}{{\mathbf{F}}}
\newcommand{\vG}{{\mathbf{G}}}
\newcommand{\vH}{{\mathbf{H}}}
\newcommand{\vI}{{\mathbf{I}}}
\newcommand{\vJ}{{\mathbf{J}}}
\newcommand{\vK}{{\mathbf{K}}}
\newcommand{\vL}{{\mathbf{L}}}
\newcommand{\vM}{{\mathbf{M}}}
\newcommand{\vN}{{\mathbf{N}}}
\newcommand{\vO}{{\mathbf{O}}}
\newcommand{\vP}{{\mathbf{P}}}
\newcommand{\vQ}{{\mathbf{Q}}}
\newcommand{\vR}{{\mathbf{R}}}
\newcommand{\vS}{{\mathbf{S}}}
\newcommand{\vT}{{\mathbf{T}}}
\newcommand{\vU}{{\mathbf{U}}}
\newcommand{\vV}{{\mathbf{V}}}
\newcommand{\vW}{{\mathbf{W}}}
\newcommand{\vX}{{\mathbf{X}}}
\newcommand{\vY}{{\mathbf{Y}}}
\newcommand{\vZ}{{\mathbf{Z}}}

\newcommand{\cA}{{\mathcal{A}}}
\newcommand{\cB}{{\mathcal{B}}}
\newcommand{\cC}{{\mathcal{C}}}
\newcommand{\cD}{{\mathcal{D}}}
\newcommand{\cE}{{\mathcal{E}}}
\newcommand{\cF}{{\mathcal{F}}}
\newcommand{\cG}{{\mathcal{G}}}
\newcommand{\cH}{{\mathcal{H}}}
\newcommand{\cI}{{\mathcal{I}}}
\newcommand{\cJ}{{\mathcal{J}}}
\newcommand{\cK}{{\mathcal{K}}}
\newcommand{\cL}{{\mathcal{L}}}
\newcommand{\cM}{{\mathcal{M}}}
\newcommand{\cN}{{\mathcal{N}}}
\newcommand{\cO}{{\mathcal{O}}}
\newcommand{\cP}{{\mathcal{P}}}
\newcommand{\cQ}{{\mathcal{Q}}}
\newcommand{\cR}{{\mathcal{R}}}
\newcommand{\cS}{{\mathcal{S}}}
\newcommand{\cT}{{\mathcal{T}}}
\newcommand{\cU}{{\mathcal{U}}}
\newcommand{\cV}{{\mathcal{V}}}
\newcommand{\cW}{{\mathcal{W}}}
\newcommand{\cX}{{\mathcal{X}}}
\newcommand{\cY}{{\mathcal{Y}}}
\newcommand{\cZ}{{\mathcal{Z}}}

\newcommand{\ri}{{\mathrm{i}}}
\newcommand{\rr}{{\mathrm{r}}}

\newcommand{\EE}{{\mathbb{E}}}


%% macros for math notions and operators

\newcommand{\RR}{\mathbb{R}}
\newcommand{\CC}{\mathbb{C}}
\newcommand{\ZZ}{\mathbb{Z}}
\renewcommand{\SS}{{\mathbb{S}}}
\newcommand{\SSp}{\mathbb{S}_{+}}
\newcommand{\SSpp}{\mathbb{S}_{++}}
\newcommand{\sign}{\mathrm{sign}}
\newcommand{\vzero}{\mathbf{0}}
\newcommand{\vone}{{\mathbf{1}}}

\newcommand{\st}{{\text{s.t.}}} % subject to
\newcommand{\St}{{\mathrm{subject~to}}} % subject to
\newcommand{\op}{{\mathrm{op}}} % subscript for operator norm
\newcommand{\opt}{{\mathrm{opt}}} % subscript for optimal solution
%\newcommand{\supp}{{\mathrm{supp}}} % support
\newcommand{\Prob}{{\mathrm{Prob}}} % probability
\newcommand{\Diag}{{\mathrm{Diag}}} % vector -> diagonal matrix
%\newcommand{\diag}{{\mathrm{diag}}} % matrix diagonal -> vector
\newcommand{\dom}{{\mathrm{dom}}} % domain
\newcommand{\range}{{\mathrm{range}}} % domain
%\newcommand{\grad}{{\nabla}}    % gradient
\newcommand{\tr}{{\mathrm{tr}}} % trace
\newcommand{\TV}{{\mathrm{TV}}} % total variation
\newcommand{\Proj}{{\mathrm{Proj}}}
\newcommand{\prj}{{\mathrm{prj}}}
\newcommand{\prox}{\mathbf{prox}}
\newcommand{\refl}{\mathbf{refl}}
\newcommand{\reflh}{\refl^{\bH}}
\newcommand{\proxh}{\prox^{\bH}}
\newcommand{\minimize}{\text{minimize}}
\newcommand{\bgamma}{\boldsymbol{\gamma}}
\newcommand{\bsigma}{\boldsymbol{\sigma}}
\newcommand{\bomega}{\boldsymbol{\omega}}
\newcommand{\blambda}{\boldsymbol{\lambda}}
\newcommand{\bH}{\vH}
\newcommand{\bbH}{\mathbb{H}}
\newcommand{\bB}{\boldsymbol{\cB}}
\newcommand{\Tau}{\mathrm{T}}
\newcommand{\tnabla}{\widetilde{\nabla}}
\newcommand{\TDRS}{T_{\mathrm{DRS}}}
\newcommand{\TPRS}{T_{\mathrm{PRS}}}
\newcommand{\TFBS}{T_{\mathrm{FBS}}}
\newcommand{\best}{\mathrm{best}}
\newcommand{\kbest}{k_{\best}}
\newcommand{\diff}{\mathrm{diff}}
\newcommand{\barx}{\bar{x}}
\newcommand{\xgbar}{\bar{x}_g}
\newcommand{\xfbar}{\bar{x}_f}
\newcommand{\hatxi}{\hat{\xi}}
\newcommand{\xg}{x_g}
\newcommand{\xf}{x_f}
\newcommand{\du}{\mathrm{d}u}
\newcommand{\dy}{\mathrm{d}y}
\newcommand{\kconvergence}{\stackrel{k \rightarrow \infty}{\rightarrow }}
\DeclareMathOperator{\shrink}{shrink} % shrinkage
\DeclareMathOperator*{\argmin}{arg\,min}
\DeclareMathOperator*{\argmax}{arg\,max}
\DeclareMathOperator*{\Min}{minimize}
\DeclareMathOperator*{\Max}{maximize}
\DeclareMathOperator*{\Fix}{Fix}
\DeclareMathOperator*{\zer}{zer}
\DeclareMathOperator*{\nablah}{\nabla^{\bH}}
\DeclareMathOperator*{\gra}{gra}
\DeclarePairedDelimiter{\dotpb}{\langle}{\rangle_{\bH}}
\DeclarePairedDelimiter{\dotpv}{\langle}{\rangle_{\vH}}
\DeclarePairedDelimiter{\dotp}{\langle}{\rangle}

%% macros for environments math equations

\newcommand{\MyFigure}[1]{../fig/#1}

\newcommand{\bc}{\begin{center}}
\newcommand{\ec}{\end{center}}

\newcommand{\bdm}{\begin{displaymath}}
\newcommand{\edm}{\end{displaymath}}

\newcommand{\beq}{\begin{equation}}
\newcommand{\eeq}{\end{equation}}

\newcommand{\bfl}{\begin{flushleft}}
\newcommand{\efl}{\end{flushleft}}

\newcommand{\bt}{\begin{tabbing}}
\newcommand{\et}{\end{tabbing}}

\newcommand{\beqn}{\begin{align}}
\newcommand{\eeqn}{\end{align}}

\newcommand{\beqs}{\begin{align*}} % no equation numbers
\newcommand{\eeqs}{\end{align*}}  % no equation numbers

%% macros for theorem-like environments

%\newtheorem{theorem}{Theorem}
\newtheorem{assumption}{Assumption}
%\newtheorem{condition}{Condition}
%\newtheorem{rul}{Rule}
%\newtheorem{definition}{Definition}
%\newtheorem{corollary}{Corollary}
%\newtheorem{remark}{Remark}
%\newtheorem{lemma}{Lemma}
%\newtheorem{proposition}{Proposition}
%\newtheorem{example}{Example}
%\newtheorem{proof}{Proof}

%%%%%%%%%%%%%%%%%%%%%%%%%%%%%%%%%%%%%%%%%%%%%

\title{Assignment 1 of \textbf{MATP4820}}
\author{Izaak King}
\date{February 2, 2026}



\begin{document}

\maketitle

\section*{Problem 1}
Let the matrix $A=\left[\begin{array}{rrr} 2 & -1 & 2 \\ 0 & -3 & 3 \end{array}\right]$. Compute $\|A\|_1$, $\|A\|_2$, and $\|A\|_\infty$.

\vspace{9mm}

\textbf{matrix 1-norm:} $||A||_{1} = \text{max}_{j}\sum_{i}|x_{ij}| = \text{max}\{(|2| + |0|), (|-1| + |-3|), (|2| + |3|)\} = \bm{5}$

\vspace{9mm}

\textbf{matrix 2-norm:}

\vspace{3mm}

First computing $AA^{T}$: 

\[AA^{T} =\left[\begin{array}{rrr} 2 & -1 & 2 \\ 0 & -3 & 3 \end{array}\right] \left[\begin{array}{rrr} 2 & 0\\ -1 & -3 \\ 2 & 3 \end{array}\right] = \left[\begin{array}{rrr} 2(2) - 1(-1) + 2(2) & 2(0) - 1(-3) + 2(3) \\ 0(2) - 3(-1) + 3(2) & 0(0) - 3(-3) + 3(3) \end{array}\right] = \left[\begin{array}{rrr} 9 & 9 \\ 9 & 18 \end{array}\right]\]

\vspace{3mm}

Now computing eigenvalues of $AA^{T}$:

\vspace{3mm}

\[ \text{det}(\lambda I - AA^{T}) = \text{det}\left[\begin{array}{ccc} \lambda - 9 & 9 \\ 9 & \lambda - 18 \end{array}\right] = (\lambda - 9)(\lambda - 18) - 81  = 0 \]

\vspace{3mm}

Solving for $\lambda:$ 

\vspace{3mm}

\[ \lambda^{2} - 27 \lambda + 81 = 0 \implies \lambda = \frac{27 \pm \sqrt{27^2 - 4(81)}}{2}\]

\vspace{3mm}

Therefore, taking largest eigenvalue, $||A||_{2} = \sqrt{\frac{27 + \sqrt{405}}{2}} = \bm{\sqrt{\frac{27 + 9 \sqrt{5}}{2}}}$

\vspace{9mm}

\textbf{matrix $\infty$-norm:} $||A||_{\infty} = \text{max}_{i}\sum_{j}|x_{ij}| = \text{max}\{(|2| + |-1| + |2|), (|0| + |-3| + |3|)\} = \bm{6}$

\newpage

\section*{Problem 2}

Compute the gradient and Hessian matrix of the function $f(x_1, x_2) = 25(x_2-x_1^2)^2 + (1-x_1)^2$. Find all local minimizers of this function. Justify your answer.

\vspace{9mm}

First expanding $f(x_{1}, x_{2})$:

\[f(x_{1},x_{2}) = 25(x_{2}^{2} - 2x_{1}^{2}x_{2} + x_{1}^{4}) + (1 - 2x_{1} + x_{1}^{2}) = 25x_{1}^{4} + 25x_{2}^{2} - 50x_{1}^{2}x_{2} + x_{1}^{2} - 2x_{1} + 1\]

\vspace{3mm}

Compute the gradient:

\[ \nabla f(x_{1}, x_{2}) = \left[\begin{array}{ccc} 100x_{1}^{3} - 100x_{1}x_{2} + 2x_{1}- 2 \\ 50x_{2} - 50x_{1}^{2} \end{array}\right] \]

\vspace{3mm}

Compute the Hessian Matrix:

\[ \nabla^{2}f(x_{1}, x_{2}) = \left[\begin{array}{ccc} 300x_{1}^{2} - 100x_{2} + 2 & - 100x_{1} \\ -100x_{1} & 50 \end{array}\right]\]

\vspace{3mm}

Now solving $\nablaf(x_{1}, x_{2}) = 0$:

\[\begin{cases}
100x_{1}^{3} - 100x_{1}x_{2} + 2x_{1}- 2  = 0\\ 
50x_{2} - 50x_{1}^{2} = 0
\end{cases}\]

\vspace{3mm}

Solving for $x_1, x_2$

\[  \implies x_{1} = 1 \implies x_{2} = 1 \]

\vspace{3mm}

So the solution is $x_{1} = 1, x_{2} = 1$, and this point is a stationary point.

\vspace{3mm}
\newpage
Now checking the Hessian Matrix at stationary point (1,1):

\[ \nabla^{2}f(1, 1) = \left[\begin{array}{ccc} 202 & - 100 \\ -100 & 50 \end{array}\right]\]

\[ \lambda_{1} + \lambda_{2} = 202 + 50 = 252 \]

\[ \lambda_{1} \lambda_{2} = 202(50) - (-100(-100)) = 100 \]

\[ \implies \lambda_{1} > 0, \lambda_{2} > 0\]

\vspace{3mm}

\noindent Thus the Hessian Matrix is positive definite, so $x_{1} = 1, x_{2} = 1,$ or (1,1) is the a local minimizer, and the only local minimizer of this function.

\newpage

\section*{Problem 3} 

Let $f(x_1, x_2) = \frac{1}{2}x_1^2 + cx_1x_2 + x_2^2 - x_1 + 2x_2 - 4$, where $c$ is a real number. Give all values of $c$ such that $f$ is convex. 

\vspace{9mm}

\noindent We will use the known fact that $f(x_{1}, x_{2}$ is a twice differentiable function, and that if $\nabla^{2}f(x)$ is Positive Semi-definite at any $x$, then $f$ is convex.

\vspace{3mm}

First, computing the gradient:

\[ \nabla f(x_{1}, x_{2}) = \left[\begin{array}{r} x_{1} + cx_{2} - 1 \\ cx_{1} + 2x_{2} + 2 \end{array}\right] \]

\vspace{3mm}

Now Hessian:

\[\nabla^{2}f(x_{1}, x_{2}) = \left[\begin{array}{rr} 1 & c \\ c & 2 \end{array}\right]\]

\vspace{3mm} 

Since $ 1 \geq 0$ and $2 \geq 0$, we just need $1(2) - c^{2} \geq 0$:

\[ 2 - c^{2} \geq 0 \implies 2 \geq c^{2} \implies -\sqrt{2} \leq c \leq \sqrt{2} \]

\vspace{3mm}

\noindent So, with $-\sqrt{2} \leq c \leq \sqrt{2}$ our Hesian matrix will be Positive Semi-definite. With the Hessian Positive Semidefinite, $f$ will be a convex function. Thus, $f$ is a convex function for $-\sqrt{2} \leq c \leq \sqrt{2}$.

\newpage

\section*{Problem 4} 

Is the sequence $x^{(k)}=1+(-\frac{3}{5})^{2^k}$ for $k=1,2,\ldots$, Q-quadratically convergent to a finite real number? If yes, give the limit and prove it; if not, explain why.

\vspace{9mm}
\noindent Yes, the sequence is Q-quadratically convergent to the finite limit $x^{*} = 1$
\vspace{9mm}

\hangindent=2em
\hangafter=1
\noindent Proof. Suppose $x^{(k)} \rightarrow x^{*}$ as $k \rightarrow \infty$. We say the convergence is Q-quadratic if there is $M > 0$ such that 

\[ \frac{||x^{(k+1)} - x^{*}||_{2}}{||x^{(k)} - x^{*}||_{2}^{2}} \leq M, 
\text{ where k is large enough}\]

\vspace{3mm}

First, we can find $x^{*}$:

\[ (-\frac{3}{5})^{2^k} \rightarrow 0 \text{ as } k \rightarrow \infty\]
\vspace{3mm}

So,

\[x^{(k)} \rightarrow 1 \text{ as } k \rightarrow \infty. \text{ Take } x^{*} = 1.\]

\vspace{3mm}
Now,



\[
\frac{\|x^{(k+1)} - x^*\|_2}{\|x^{(k)} - x^*\|_2^2}
= \frac{\left|1 + (-\frac{3}{5})^{2^{k+1}} - 1 \right|}{\left|1 + (-\frac{3}{5})^{2^k} - 1\right|^2}
= \frac{\left(\frac{3}{5}\right)^{2^{k+1}}}{\left(\frac{3}{5}\right)^{2\cdot 2^k}}
= \frac{\left(\frac{3}{5}\right)^{2\cdot 2^k}}{\left(\frac{3}{5}\right)^{2\cdot 2^k}}
= 1.
\]




\vspace{3mm}

Thus,

\[ \frac{||x^{(k+1)} - x^{*}||_{2}}{||x^{(k)} - x^{*}||_{2}^{2}} \leq M, \text{ with } M = 1 > 0. \]

\vspace{3mm}

Therefore, we have shown that the sequence is Q-quadratically convergent to $x^{*} = 1 $

\begin{flushright}
$\square$
\end{flushright}

\end{document}
